\newpage
\appendix

\section{Motivation Study Setup}
\label{sec:appendix:motivation-setup}

% Sources: ../FedLLM/exp/_archive_exps/1023_flan1_motivation/PROTOCOL.md and
% ../FedLLM/exp/_archive_exps/1023_flan1_motivation/recovery_tail_01/report/flan8c/RESULTS.md.
% Frozen configurations: ../FedLLM/exp/_archive_exps/1023_flan1_motivation/configs/
% flan1_8c_two_ranks_dir05_q8.yaml and flan1_8c_two_ranks_dir05_geometry_q8.yaml.

This appendix details the shared-width study and the training-only similarity diagnostic in \figref{fig:acl-motivation}.

\fakeparagraph{Dataset and Clients}
The studies use eight task categories from FLAN~\cite{wei2022finetuned}: coreference, entailment, linguistic acceptability, paraphrase, question classification, structure-to-text generation, text formatting, and word disambiguation.
We use a Dirichlet $\alpha=0.5$ allocation across eight clients, balanced to give each client 300 training examples and 200 test examples with matching local train/test task proportions.
Four clients have fixed total rank 16, and the other four have rank 64.
The dataset and client rank assignments are shared by both studies and remain fixed across the width sweep.

\fakeparagraph{Model and Training}
Both studies use Qwen3.5-0.8B, seed 615, full client participation, and 10 communication rounds with 10 local optimizer steps per round.
The microbatch size is 2 with gradient accumulation over 4 microbatches, and the initial learning rate is $10^{-4}$, decayed by a factor of 0.99 per round.
The physical LoRA rank and scaling parameter are both 64, with active capacity fixed to each client's assigned rank and dropout 0.05.
The maximum sequence length is 2,048 tokens, reserving up to 512 tokens for the answer.
We use the base shared-prefix method: for shared width $s>0$, the $B$ columns and $A$ rows of the first $s$ LoRA components are averaged separately across all clients, including the receiving client, after each round, while the remaining active components stay local.
Aggregation is weighted by training sample count, giving equal weights here because all clients have 300 training examples.
Both portions participate in training and prediction; $s=0$ uses the same local training protocol with no shared components.
The same factor-averaging rule is used at every nonzero width, with no partner selection or fitted reception strength.

\fakeparagraph{Shared-Width Study}
For \figref{fig:acl-motivation}(a), all clients use a common shared width $s$, swept over $s\in\{0,\ldots,8\}$ while keeping the model, data, client ranks, seed, and training budget fixed.
We evaluate the final post-aggregation models at R9 on all 200 test examples per client and average ROUGE-L first within each client, then equally across the four clients in each rank group.
The plotted change is each group's mean at $s$ minus its own mean at $s=0$, in percentage points, using the full-precision saved client scores.
The rank-16 and rank-64 groups have $s=0$ ROUGE-L scores of $64.2584\%$ and $65.3399\%$, respectively.
The overall client means at $s=0$ and $s=8$ are $64.7991\%$ and $64.7644\%$, respectively.
The two rank groups contribute equally to the overall mean because each contains four clients.

\fakeparagraph{Similarity Diagnostic}
For \figref{fig:acl-motivation}(b), all clients use a fixed shared width $s=8$ in a separate training-only rerun, with a snapshot after every round and no test evaluation.
Its scientific configuration matches the corresponding width-sweep arm except for saving all round snapshots instead of only the final one.
R0 is the first completed training round, and R9 is the tenth and final round.
The diagnostic is computed separately for the total-rank groups 16 and 64.
Cosines compare the effective shared uploads $B_{\mathrm{upload}}A_{\mathrm{upload}}-B_{\mathrm{global}}A_{\mathrm{global}}$ relative to the pre-round global state, concatenated across all LoRA layers.
For each total rank $r$, the same-rank mean uses different-domain pairs with both clients at $r$; the same-domain mean uses same-domain pairs with different ranks and one client at $r$.
Domain denotes the recorded dominant task category.
Each rank group has six same-rank/different-domain pairs and one same-domain/different-rank pair, with membership fixed across rounds.
The same-domain reference is a single cross-rank pair with entailment as its dominant category, and both curves use this reference.
All selected cosines are defined in every round; the complete diagnostic contains 28 unordered pairs per round and 280 pair/round records in total.
Each curve plots the same-rank pair mean minus the same-domain pair mean, so positive values indicate higher same-rank cosine.

\fakeparagraph{Interpretation}
Both studies use one seed and provide no across-seed uncertainty estimate.
Rank groups differ in task composition, so their contrasts are descriptive rather than causal rank effects.
The matched $s=0$ arm permits within-protocol comparisons, but the test curves are descriptive and are not used to select shared widths.
Changing $s$ jointly changes protected capacity, exchanged content, and training, so these comparisons do not isolate donor-specific transfer or establish a general advantage for sharing.
The shared same-domain reference is sparse, and the two similarity curves are not independent repeats.
The two panels use separate runs and endpoints; the cosine diagnostic does not measure benefit to the target client or validate a partner-selection rule.
Because the averaging rule is fixed in the shared-width sweep, the two panels do not measure an allocation-by-collaboration interaction.

\fakeparagraph{Relation to the Proposed Method}
The probe couples reception and contribution through the same global width $s$ and averages the shared factors directly, following the selective rank sharing definition in \secref{sec:problem}.
In \secref{sec:method}, clients instead contribute their complete adapters, while $s_i$ limits the receiving subspace for a projected effective update.
Consequently, the probe supplies observations for the research questions, while evaluating the proposed reception rule requires its own matched controls.

\section{MMLU-Pro Baseline Results}
\label{sec:appendix:mmlu-baselines}

% MMLU-Pro is kept here pending the user's decision on its inclusion in the final paper.
\tabref{tab:exp-mmlu-dir05} reports client-macro multiple-choice (MC) accuracy on MMLU-Pro under Dirichlet $\alpha=0.5$.
The experiments use Qwen3.5-0.8B, seed 615, and the final checkpoint after 10 communication rounds.
The 13 clients have total rank budgets 16, 32, and 64, assigned to five, five, and three clients, respectively.
Accuracy is computed on each client's full original test set and averaged equally across clients.

\begin{table}[t]
\centering
\small
\setlength{\tabcolsep}{4pt}
\caption{Client-macro baseline MC accuracy (\%) on MMLU-Pro under Dirichlet $\alpha=0.5$.
The highest and second-highest scores are shown in bold and underlined, respectively.}
\label{tab:exp-mmlu-dir05}
\begin{tabular}{lr}
\toprule
\textbf{Method} & \textbf{MC accuracy (\%)} \\
\midrule
LocalOnly                          & 30.0322 \\
FlexLoRA                           & \textbf{31.9374} \\
FLoRA                              & \underline{30.5800} \\
\bottomrule
\end{tabular}
\end{table}

FlexLoRA has the highest MC accuracy among the reported baselines.
